%% kapa: exam=1 round=26 year=2015 period=2 form=A subject=appraisal_law subject_ko=부동산관계법규 seq=21 number=21
%% source: Q-Net 감정평가사 기출문제 (cst00309) artlSeq=5138001 file=extracted/1008/5138001/제26회_감정평가사 1차_기출/제26회_감정평가사 1차_기출/2교시_감평_A형.pdf
%% conversion: kordoc 4.24.51 -> tools/gen1.py (auto); review=proofread-sonnet; flags=none
\begin{question}{21}
부동산가격공시 및 감정평가에 관한 법령상 용어에 관한 설명으로 옳지 않은 것은?
\begin{choices}
  \item ‘표준지공시지가’라 함은 「부동산가격공시 및 감정평가에 관한 법률」의 규정에 의한 절차에 따라 국토교통부장관이 조사ㆍ평가하여 공시한 표준지의 단위면적당 가격을 말한다.
  \item ‘토지등’에는 「입목에 관한 법률」에 의한 입목은 포함되지 아니한다.
  \item ‘적정가격’이라 함은 당해 토지 및 주택에 대하여 통상적인 시장에서 정상적인 거래가 이루어지는 경우 성립될 가능성이 가장 높다고 인정되는 가격을 말한다.
  \item ‘감정평가’라 함은 토지등의 경제적 가치를 판정하여 그 결과를 가액으로 표시하는 것을 말한다.
  \item ‘다세대주택’이라 함은 주택으로 쓰이는 1개 동의 연면적(지하주차장 면적을 제외한다)이 660제곱미터 이하이고, 층수가 4개층 이하인 주택을 말한다.
\end{choices}
\end{question}
\answer{2}
